\begin{workedexample}[Worked Example: Chemical-shift displacement]
Fat and water precess about $3.5\ \text{ppm}$ apart. At $1.5\ \text{T}$
($64\ \text{MHz}$) that is $3.5\times10^{-6}\times64\times10^6 = 224\ \text{Hz}$; at
$3\ \text{T}$ it doubles to $448\ \text{Hz}$. Because the readout maps frequency to
position, fat shifts spatially relative to water --- worse at high field and with
low readout bandwidth.
\end{workedexample}

\begin{keytakeaways}
\begin{itemize}[leftmargin=1.4em,itemsep=2pt]
\item Chemical shift, susceptibility, motion, aliasing, and Gibbs ringing are the common artefacts.
\item Chemical-shift and susceptibility effects scale with $B_0$ (worse at high field).
\item Higher readout bandwidth reduces chemical-shift displacement; oversampling cures aliasing.
\end{itemize}
\end{keytakeaways}

\begin{exercises}
\item Fat--water shift (Hz) at $3\ \text{T}$ ($3.5\ \text{ppm}$)?
\item Does the susceptibility artefact improve or worsen at $7\ \text{T}$?
\item What reduces chemical-shift displacement in the readout?
\end{exercises}

\furtherreading{Haacke~\cite{haacke} (ch.~20); Nishimura~\cite{nishimura}.}
